\documentclass[12pt,a4paper]{article}
\usepackage{graphicx}
\usepackage{amsmath}
\usepackage{hyperref}
\usepackage{geometry}
\geometry{margin=2.5cm}
\emergencystretch=3em

\title{\textbf{CDFD Runtime Paper 12: Candidate Laws, Discovery Program, and Responsible Use}}
\author{Steve Bico Mujjabi, MD\\
\small Independent Researcher\\
\small Founder, Vura Labs\\
\small Kampala, Uganda\\
\small ORCID: \href{https://orcid.org/0009-0001-0556-5516}{0009-0001-0556-5516}}
\date{May 2026}

\begin{document}
\maketitle

\begin{abstract}
\noindent \textbf{Executive synthesis.} Earlier drafts called the final synthesis
``immutable laws.'' This upgraded paper uses more disciplined language:
candidate laws, model invariants, and falsification targets. CDFD Runtime
experiments suggest recurring motifs--adaptive constraint, operating-ratio
limits, structural memory, weakest-link throughput, and Life Number thresholds.
These motifs are useful research guides, not final proof that reality has been
fully decoded.
\end{abstract}


\paragraph{Greek-symbol convention.}
Unless a manuscript states a tighter equation-local meaning, Greek symbols are local CDFD variables or coefficients, not universal constants. Here $\Phi$ denotes driving flux, $\Psi_s$ denotes the adaptive operating ratio, $\Omega$ denotes a domain or state space, and $\Lambda$ denotes the Life Number where used. The coefficients $\alpha$, $\beta$, and $\gamma$ denote accumulation or reinforcement, relaxation or decay, and diffusion, coupling, or redistribution. The symbols $\delta$ and $\epsilon$ denote perturbation or tolerance scales, while $\eta$, $\lambda$, $\mu$, $\nu$, $\rho$, $\sigma$, $\tau$, and $\phi$ denote local efficiency, rate, baseline, frequency, density or correlation, noise or variance, time-scale, and phase or field terms. Subscripts restrict any coefficient to the named pathway, tissue, domain, or experiment.

\begin{figure}[ht]
\centering
\fbox{\begin{minipage}{0.92\textwidth}
\centering
\textbf{Runtime evidence path}\\[0.4em]
Prior CDFD mathematics $\longrightarrow$ CDFL/runtime state $\longrightarrow$ bounded output $\longrightarrow$ falsification target.
\end{minipage}}
\caption{Conceptual schematic for the runtime papers. The engine translates the mathematical frame from the prior CDFD papers into executable CDFL/runtime diagnostics; candidate outputs remain tied to explicit tests before any stronger claim is made.}
\end{figure}

\section{Prior CDFD Basis}
This manuscript does not rederive the mathematical core of CDFD. It uses the
Mujjabi Adaptive Operating Ratio and the constrained-vacuum notation introduced
in Part I \cite{MujjabiI2026}, the torus-knot hierarchy and boundary-mode work
from the Part I sequence \cite{MujjabiVI2026,MujjabiIX2026}, the public balance
law developed in the Part I capstone \cite{MujjabiX2026}, the Part I transport
tests \cite{MujjabiXII2026}, and the Life Number and tri-regime biology bridge
from Part II \cite{MujjabiPartII2026}. The runtime papers therefore act as an
execution layer: they keep the mathematics connected to commands, outputs,
falsification tests, and domain-specific limits.


\section{Why the Name Changes}
The phrase ``immutable laws'' overstates the current evidence. The runtime has
scripts, diagnostics, and reports. It does not yet have enough independent
empirical validation to call every recurrent pattern a universal law. The final
paper therefore names a research program: candidate Mujjabi laws that remain open
to testing. Named CLI-facing tests are listed in
\texttt{MUJJABI\_RUNTIME\_LAWS\_AND\_TESTS.md}, including aromatic source-mix and
photochemical endpoint guardrails added in the May 2026 runtime pass.

\section{Candidate Mujjabi Law 1: Adaptive Constraint}
\textit{Sustained flow tends to generate, reveal, or amplify constraint.}
In runtime form, this appears through the constraint equation
\begin{equation}
    \partial_t C = \alpha|\partial_t \Phi|-\beta C+\gamma\nabla^2 C.
\end{equation}
The falsification target is simple: find mapped systems where increasing drive
does not produce any measurable constraint response and where that absence is
not already explained by a saturated or buffered medium.

\section{Candidate Mujjabi Law 2: Operating-Ratio Regimes}
\textit{The relation between $\Phi$ and $C$ predicts regime transitions.}
The operating ratio
\begin{equation}
    \Psi_s=\frac{\Phi}{C}S M_s
\end{equation}
is the central runtime diagnostic. It must be validated separately in each
domain. A high $\Psi_s$ in one domain may mean overload; in another it may mean
transient productive growth. Interpretation is not automatic.

\section{Candidate Mujjabi Law 3: Structural Memory}
\textit{Overload can leave retained memory in the medium.}
The vacuum-hysteresis experiment selected in
\texttt{FINAL\_RELEASE\_EXPERIMENT\_SELECTION.md} records a localized $M_s$
residual. This supports the runtime's memory mechanism as a candidate diagnostic
and a measurement protocol. It does not prove a physical vacuum-memory field.

\section{Candidate Mujjabi Law 4: Weakest-Link Throughput}
\textit{Network throughput is often limited by its most constrained effective
channel.}
This appears in origins-of-life, medicine, networks, infrastructure, and supply
chains. The generic diagnostic is useful, but real systems require real capacity
data. The law fails where redundancy, storage, or rerouting invalidates the
simple bottleneck approximation.

\section{Candidate Mujjabi Law 5: Life Number Threshold}
\textit{Sustained life-like throughput requires energy input, electron
transport, proton coupling, relaxation, and stabilization to balance.}
The Life Number is
\begin{equation}
    \Lambda =
    \frac{\Phi \sigma_e \sigma_p \tau_{\mathrm{relax}}}
    {S_{\mathrm{stab}}E_{\mathrm{maintenance}}}.
\end{equation}
The OOL phase experiment treats $\Lambda=1$ as a candidate transition boundary.
The claim remains model-bound until laboratory and geochemical tests challenge
the mapping.

\section{Discovery Program}
The local discovery program now has four evidence levels:
\begin{enumerate}
    \item selected outputs such as the Mujjabi Vacuum Hysteresis Test, the
    Mujjabi \(n=7\) Knot Stability Test, and the Mujjabi Life Number Boundary
    Test;
    \item broad sweeps such as frontier and gigamarathon outputs;
    \item reports that define candidate interpretations and falsification
    conditions;
    \item CLI and CDFL commands that make future runs reproducible.
\end{enumerate}

\section{Responsible Use}
No runtime output is medical advice, device certification,
financial policy, or physical proof without external validation. CDFD is a
framework for asking sharper questions. Its power increases when its claims are
clearly bounded.

\section{Conclusion}
The synthesis of the runtime series is not certainty. It is a disciplined
research engine: define variables, run the model, save the output, expose the
failure condition, and let better evidence decide which candidate laws survive.
\begin{thebibliography}{9}
\bibitem{MujjabiI2026}
Steve Bico Mujjabi, MD. \emph{Vortex Stability in a Constrained Transport Vacuum and the Origin of the Fine-Structure Constant}. CDFD Part I Paper I, preprint, 2026.

\bibitem{MujjabiVI2026}
Steve Bico Mujjabi, MD. \emph{The Universal Torus Knot Hierarchy: $Z_n$ Symmetry, the Cinquefoil Family, and the General Structure of CDFT Particle Families}. CDFD Part I Paper VI, preprint, 2026.

\bibitem{MujjabiIX2026}
Steve Bico Mujjabi, MD. \emph{Even-$n$ Torus Knots in CDFT: The Exclusion of $T(2,2)$, the Extension to $T(2,2m)$, and the Anti-Phase Mass Scaling Law}. CDFD Part I Paper IX, preprint, 2026.

\bibitem{MujjabiX2026}
Steve Bico Mujjabi, MD. \emph{The Vacuum Equation of State: Deriving Torus Knot Self-Consistency from the Public CDFD Balance Law $\Psi = \Phi/C$}. CDFD Part I Paper X, preprint, 2026.

\bibitem{MujjabiXII2026}
Steve Bico Mujjabi, MD. \emph{Friction, Superconductivity, Plasma States, and Transport Mysteries: Observable Tests of Adaptive Constraint}. CDFD Part I Paper XII, preprint, 2026.

\bibitem{MujjabiPartII2026}
Steve Bico Mujjabi, MD. \emph{CDFD Part II: Origins of Life and Tri-Regime Bioenergetics}. Public release archive, 2026, \url{https://doi.org/10.5281/zenodo.20264779}.
\end{thebibliography}
\end{document}
